\chapter{Nature}
\label{ch:nature}

\section{An accident worth keeping}

Cuba has the least degraded large-island ecology in the Caribbean, and it has it
substantially by accident.

The country retains extensive coral reef systems, including at the Jardines de la
Reina --- an archipelago off the southern coast, closed to commercial fishing
since 1996, where the shark and grouper populations are among the last in the
region that resemble what Caribbean reefs looked like before industrial fishing.
It retains the Ci\'enaga de Zapata, the largest wetland in the insular Caribbean,
with a bird list including several endemics found nowhere else. It retains
Alejandro de Humboldt National Park in the east, one of the most biologically
diverse island sites on earth by area, and the mogotes of Vi\~nales. Cuban
endemism is high across plants, birds, reptiles and molluscs, as it is on any
large old island, and unusually much of it is still there.

The reasons are not virtuous. Cuba stopped being able to afford fertiliser and
pesticide in 1990. It stopped being able to build coastal hotels at Caribbean
rates. It never developed the cruise-and-condominium coastline that has consumed
much of the region. Its citizens could not buy boats or second homes. Poverty is
an effective, brutal and temporary conservation policy, and every year of the
recovery this book proposes will make it less effective.

There is also a real institutional element, which should be credited. Cuba
protected the Jardines de la Reina deliberately and enforced it. It has a
national system of protected areas covering a substantial share of territory, a
serious environmental law from 1997, a genuine tradition of field biology, and in
Tarea Vida a national climate adaptation programme adopted in 2017 that is
scientifically better than most in the region and almost entirely unfunded. The
knowledge exists. The money and the enforcement capacity do not.

\section{The mechanism that destroys it}

The pattern is the same everywhere in the Caribbean and it is worth stating as a
mechanism rather than as a warning, because the mechanism is what a design can
interrupt.

A country decides to grow tourism. Beachfront becomes the most valuable land, so
it is built on --- and building on it removes the dune, the mangrove and the sea
grass, which is what was holding the sand and filtering the runoff. The
construction and the new population produce sewage that the treatment
infrastructure does not handle, so nutrients enter the water. Nutrients feed
algae; algae smother coral. Fishing intensifies to feed the visitors, removing
the herbivorous fish that graze the algae off the coral. Boats anchor on the
reef. Roads and causeways cut water circulation in the lagoons behind the cays.
Within fifteen to twenty-five years the reef that the destination was sold on is
dead, the beach it protected is eroding, and the country is defending its
coastline with imported rock.

Cuba is at the beginning of this sequence rather than the end, and it has one
enormous advantage: it can see the whole curve, in detail, in a dozen
jurisdictions within an hour's flight.

\section{The exposure}

Climate risk in Cuba is not a future scenario; it is a schedule.

\emph{Hurricanes} are the acute risk, and the trend in the intensity of the
strongest storms is upward. Cuba's civil defence is genuinely excellent and its
death tolls are among the lowest in the region for storms of comparable
strength --- this is a real and underappreciated Cuban achievement --- but the
physical damage is not preventable by evacuation. The 2008 season alone cost
something on the order of a fifth of national output.

\emph{Sea level} is the chronic one. Cuba's own Tarea Vida assessments project
the loss of low coastal settlements over the century and identify communities
that should not be rebuilt where they stand --- which is a politically
extraordinary thing for a government to publish, and which it did.

\emph{Saline intrusion} into coastal aquifers is already measurable and is a
water-supply problem for Havana and for the sugar plain.

\emph{Coral bleaching} follows sea temperature, and the Caribbean has had
repeated regional bleaching events. The Cuban reefs are not exempt; what they
have is better condition going in, which is the single largest determinant of
whether a reef recovers from a bleaching event.

\emph{Drought} in the eastern provinces --- the poorest, per
Chapter~\ref{ch:premises} --- has been severe and recurrent, and it interacts
with everything in Chapter~\ref{ch:land}.

\section{The design}

The argument throughout is economic rather than sentimental, because the
sentimental version loses to a hotel proposal every time.

The reefs, the wetlands and the cays are the physical stock from which
Chapter~\ref{ch:tourism}'s high-value products are sold. A dive day at the
Jardines de la Reina earns more per visitor than any beach room in the country
and it exists solely because somebody drew a line and enforced it. That is not a
constraint on the tourism strategy; it is the tourism strategy's inventory, and
the correct treatment of inventory is not to consume it.

\begin{design}{Protecting the stock}
\label{d:protect}
\begin{rules}
\rl{1}{An independent regulator} An environmental authority with the tenure,
budget and publication guarantees of Design~\ref{d:courts}, independent of the
tourism ministry, of the entities in Design~\ref{d:gaesa}, and of any ministry
that owns a development. Consent decisions and their reasons published; refusals
published. At present the promoter and the consenting authority are frequently
the same state.

\rl{2}{Coastal setback, without exception} A statutory construction setback from
the mean high-water line, applied to every developer including state and
military-linked ones, with no variance procedure. Every Caribbean setback rule
that has a variance procedure has failed through it, which is why this clause is
about the absence of the exception rather than the distance.

\rl{3}{No new causeways} No further \emph{pedraplenes} to the cays; existing ones
assessed for remediation with culverts to restore circulation. The damage these
have done to lagoon systems is documented and is reversible in part.

\rl{4}{Sewage before rooms} No accommodation development is consented without
treatment capacity commissioned and operating. Nutrient loading is the single
largest controllable driver of reef loss, and it is an infrastructure decision
taken years before the reef dies.

\rl{5}{Fishing} No-take zones maintained and extended around reef systems, with
enforcement funded from visitor revenue under clause~1 of
Design~\ref{d:natfinance}. The Jardines de la Reina model works and is
replicable.

\rl{6}{Environmental data published} Reef condition, water quality, protected
area extent and enforcement actions published annually, per
Design~\ref{d:minimum}.3.
\end{rules}
\end{design}

\begin{design}{Paying for it}
\label{d:natfinance}
\begin{rules}
\rl{1}{Visitor finance} A share of the accommodation levy and of protected-area
access fees is assigned by statute to the protected-area system and retained by
it, on the Design~\ref{d:heritage} model. Conservation that depends on an annual
budget negotiation loses the negotiation.

\rl{2}{Price the access} Where a site is at its carrying capacity, access is
priced and allocated by published auction or ballot
(Design~\ref{d:limits}). This converts scarcity into revenue for the thing that
is scarce.

\rl{3}{Carbon and ecosystem finance} Mangrove and forest carbon, and
payment-for-ecosystem-services arrangements, pursued as a revenue stream. These
markets are imperfect and the amounts are modest, but they are hard currency for
not doing something, which is unusually well suited to Cuba's position.

\rl{4}{Science as an export} Cuban field biology plus intact ecosystems plus
foreign research funding is a real and immediate line of business, and it
overlaps with Chapter~\ref{ch:tourism}'s high-value visitor and
Chapter~\ref{ch:talent}'s academic opening. It requires only that foreign
scientists can obtain permits without a political assessment.
\end{rules}
\end{design}

\begin{design}{Building for the century}
\label{d:adapt}
\begin{rules}
\rl{1}{Fund Tarea Vida} The existing national adaptation programme is
scientifically sound and unfunded. It is assigned a standing allocation from the
stabilisation fund of Design~\ref{d:rentrule}.2, which is exactly the kind of
capital spending that fund exists for.

\rl{2}{Building codes} Hurricane-resistant construction standards, enforced
through the municipal building control of Design~\ref{d:municipal}.3, applied to
housing as well as to hotels. Cuba's housing stock is its most storm-vulnerable
asset and its rebuilding under Chapter~\ref{ch:capital} is the once-in-a-century
opportunity to raise the standard.

\rl{3}{Managed retreat} Where Tarea Vida identifies settlements that cannot be
defended, relocation is planned, funded, consented and compensated, and no
reconstruction is publicly financed in place. This is politically painful and
cheaper than the alternative, and Cuba has already done the analytical work.

\rl{4}{Water} Aquifer protection, leak reduction in the distribution network ---
where losses are very large --- and water pricing above a free basic allowance.
Cuba currently loses a great deal of pumped water to leaks, which is
simultaneously a water problem and, since the water is pumped, an electricity
problem under Chapter~\ref{ch:energy}.

\rl{5}{Insurance and reserve} Per Design~\ref{d:storm}.5, an explicit published
reserve against catastrophe losses, since commercial cover for Cuban assets is
expensive and in several markets unavailable.
\end{rules}
\end{design}

\section{The trade-off, stated}

This chapter constrains Chapter~\ref{ch:tourism} and Chapter~\ref{ch:capital}
and it should say so rather than pretend the objectives are automatically
aligned.

Setbacks, sewage requirements, carrying capacity limits and no-take zones reduce
the number of hotel rooms that can be built and the speed at which they can be
built. In the short run that is less investment, fewer construction jobs and less
revenue, and a government under fiscal pressure will be offered a great deal of
money to make an exception --- which is why Design~\ref{d:protect}.2 has no
variance procedure and why the regulator's independence is written in the same
terms as the judiciary's.

The case for accepting the constraint is Chapter~\ref{ch:tourism}'s: Cuba is not
competing on volume, cannot win on volume, and is proposing to sell the one thing
in the Caribbean that is genuinely scarce. Every other country in the region has
already made the opposite trade and can report the result.
